% Zumdahl Chapter 15 -- Acid-Base Equilibria. Elaborated notes, 5 blocks.
% Back half of CED Unit 8 (buffers + titrations) and CED 7.12 (common ion).
\documentclass{apchem-notes}

\apcunitword{Chapter}
\apcunit{15}
\apcunittitle{Acid--Base Equilibria}
\apcdoctitle{Guided Notes}
\apcsubtitle{Zumdahl \S15.1--15.5 \textbullet\ PDF pp.~755--803 \textbullet\ 5 blocks}

\begin{document}
\apctitleblock

\begin{apcnote}[How this chapter fits the AP course]
This is the \textbf{back half of Unit 8} and the chapter that finishes it:
\S15.1 $\to$ CED 7.12, \S15.2 $\to$ 8.4 and 8.8--8.9, \S15.3 $\to$ 8.10,
\S15.4 $\to$ 8.5, \S15.5 $\to$ 8.5 (indicator choice).

\textbf{It also closes a gap left open in Chapter 13.} Unit 7's last two
topics were missing from the equilibrium chapter: \S15.1 supplies
\textbf{7.12, the common-ion effect}. Only \textbf{7.11 (Ksp)} is still
outstanding, and it lives in \S16.1.

\textbf{Marked enrichment:} \S15.6 (polyprotic titrations). You are
responsible for recognizing the \emph{shape} of a polyprotic curve and what
each plateau means; the multi-equivalence-point calculations are not on the
framework.

Nothing here is a new kind of problem. A buffer is a weak-acid equilibrium
with the conjugate base already present, and a titration is a
\emph{stoichiometry} problem followed by whichever equilibrium is left over.
The skill being built is deciding \emph{which} of those two you are in.
\end{apcnote}

% =====================================================================
\blocklesson{The Common Ion Effect and What a Buffer Is}{Zumdahl \S15.1--15.2}

\begin{objectives}
  \item Predict the effect of a common ion using Le Ch\^atelier's principle.
  \item Recognize a buffer and calculate its pH from $K_a$ and the
        concentration ratio.
\end{objectives}

\reading{Zumdahl \S15.1--15.2}{756--765}

\begin{warmup}[4]
  \item pH of \SI{0.10}{\Molar} \ce{HC2H3O2} ($K_a = 1.8\times10^{-5}$):
        \answerblank[2.4cm]{2.87}
  \item $K_b$ for \ce{C2H3O2^-}: \answerblank[3.2cm]{$5.6\times10^{-10}$}
  \item If $Q < K$, the equilibrium shifts: \answerblank[2.4cm]{right}
  \item Adding a product to a system at equilibrium shifts it:
        \answerblank[2.4cm]{left}
\end{warmup}

\segment{INSTRUCTION A}{25}{The common ion effect}

\notesectionz{Le Ch\^atelier, wearing a different hat}{15.1}
\practice{6}

A \term{common ion} is an ion that is already a participant in the
equilibrium you care about. Adding it shifts the position of that
equilibrium \answerblank[2.6cm]{backward} --- there is no new principle
here at all. This is Le Ch\^atelier's principle applied to an acid
ionization.

\[ \ce{HF(aq) <=> H+(aq) + F^-(aq)} \]

Dissolve \ce{NaF} in this solution and you have added \ce{F-}, a product.
The equilibrium shifts \answerblank[2cm]{left}, so the acid ionizes
\emph{less} and the solution is \answerblank[2.4cm]{less acidic} than the
acid alone.

\begin{workedexample}[Worked example 1: Zumdahl's fluoride comparison]
Compare \SI{1.0}{\Molar} \ce{HF} alone with a solution that is
\SI{1.0}{\Molar} in \ce{HF} \emph{and} \SI{1.0}{\Molar} in \ce{NaF}.
$K_a(\ce{HF}) = 7.2\times10^{-4}$.

\textbf{\ce{HF} alone.} The usual weak-acid ICE:
$x = \sqrt{(7.2\times10^{-4})(1.0)} = 2.7\times10^{-2}$, so
$\mathrm{pH} = \mathbf{1.57}$ and the acid is $\mathbf{2.7\%}$ ionized.

\textbf{With \ce{NaF} present.} Now $[\ce{F-}]_0 = 1.0$, not zero:
\[ K_a = \frac{x(1.0+x)}{1.0-x} \approx \frac{x(1.0)}{1.0}
   \quad\Longrightarrow\quad x = K_a\frac{[\ce{HF}]}{[\ce{F-}]}
   = 7.2\times10^{-4} \]
so $\mathrm{pH} = \mathbf{3.14}$ and the acid is only $\mathbf{0.072\%}$
ionized.

The common ion suppressed the ionization by a factor of about
\textbf{37}, and moved the pH by more than \textbf{1.5 units}.
\end{workedexample}

\begin{apctrap}
Notice what the algebra collapsed to:
\[ [\ce{H+}] = K_a\,\frac{[\ce{HA}]}{[\ce{A^-}]} \]
No square root, no quadratic. When \emph{both} members of a conjugate pair
are present in comparable amounts, the ICE table's $x$ is negligible against
both of them, and the whole problem reduces to a ratio. This single
expression is the engine of the entire chapter.
\end{apctrap}

\segment{GUIDED PRACTICE}{15}{Predicting the shift}

\begin{enumerate}[leftmargin=*,itemsep=0.35em]
  \item \ce{NaCN} is added to a solution of \ce{HCN}. The pH
        \answerblank[4.6cm]{rises (less ionization)}
  \item \ce{NH4Cl} is added to a solution of \ce{NH3}. The pH
        \answerblank[4.6cm]{falls (less \ce{OH-})}
  \item \ce{NaCl} is added to a solution of \ce{HC2H3O2}. The pH
        \answerblank[7.2cm]{no change --- neither ion participates}
  \item Which ion in item 3 would have mattered?
        \answerblank[4.6cm]{\ce{C2H3O2^-}, not \ce{Cl-}}
\end{enumerate}

\segment{INSTRUCTION B}{20}{What makes a solution a buffer}

\notesectionz{Definition and recognition}{15.2}
\practice{6}

A \term{buffered solution} resists a change in pH when acid or base is
added. It contains \textbf{significant amounts of both members of a
conjugate acid--base pair}: a weak acid together with its conjugate base, or
a weak base together with its conjugate acid.

Three ways a buffer shows up on the exam:

\begin{enumerate}[leftmargin=*,itemsep=0.3em]
  \item \textbf{Stated outright} --- \ce{HC2H3O2} plus \ce{NaC2H3O2}, or
        \ce{NH3} plus \ce{NH4Cl}.
  \item \textbf{Made by partial neutralization} --- a weak acid plus
        \emph{less than enough} strong base to consume it. The base
        converts part of the acid to its conjugate base, and what remains is
        a mixture of both.
  \item \textbf{Hiding at the halfway point of a titration}, which is item 2
        with the amounts chosen so the two are equal.
\end{enumerate}

\begin{apctrap}
A strong acid and its ``conjugate base'' is \emph{not} a buffer.
\ce{HCl} and \ce{NaCl} do nothing: \ce{Cl-} has no affinity for protons, so
there is nothing to absorb added acid. A buffer requires a \emph{weak}
conjugate pair --- the reversibility is the whole point.
\end{apctrap}

\segment{APPLICATION}{20}{Ratio calculations}

\begin{enumerate}[leftmargin=*,itemsep=0.4em]
  \item Find $[\ce{H+}]$ and pH for a solution \SI{0.10}{\Molar} in \ce{HF}
        and \SI{0.30}{\Molar} in \ce{NaF}. \workspace[2.4cm]
        \keyonly{\begin{apcanswer}$[\ce{H+}] = (7.2\times10^{-4})
        (0.10/0.30) = 2.4\times10^{-4}$; pH $= \mathbf{3.62}$\end{apcanswer}}
  \item Identify each as a buffer or not, with a reason:
        \begin{enumerate}[leftmargin=1.4em,itemsep=0.25em]
          \item \ce{HNO3} $+$ \ce{NaNO3}
                \answerblank[7cm]{no --- \ce{HNO3} is strong}
          \item \ce{NH3} $+$ \ce{NH4Cl}
                \answerblank[8cm]{yes --- weak base and its conjugate acid}
          \item \ce{HCN} $+$ \ce{NaCN}
                \answerblank[8cm]{yes --- weak acid and its conjugate base}
          \item \ce{NaOH} $+$ \ce{NaCl}
                \answerblank[8cm]{no --- strong base, no conjugate pair}
        \end{enumerate}
\end{enumerate}

\begin{exitticket}
A student says ``\ce{NaF} makes the solution basic, so adding it to \ce{HF}
raises the pH.'' The conclusion is right but the reasoning is incomplete.
Give the better explanation.
\keyonly{\begin{apcanswer}\ce{F-} is a product of the \ce{HF} ionization, so
adding it shifts that equilibrium to the left and suppresses the production
of \ce{H+}. The hydrolysis of \ce{F-} is a real but much smaller effect; the
dominant reason is the common-ion shift.\end{apcanswer}}
\end{exitticket}

% =====================================================================
\blocklesson{The Henderson--Hasselbalch Equation}{Zumdahl \S15.2}

\begin{objectives}
  \item Derive and apply $\mathrm{pH} = \mathrm{p}K_a +
        \log([\text{base}]/[\text{acid}])$.
  \item Explain why the ratio, not the absolute concentration, sets the pH.
\end{objectives}

\reading{Zumdahl \S15.2}{763--767}

\begin{warmup}[3]
  \item $[\ce{H+}] = K_a \times$ \answerblank[3cm]{$[\ce{HA}]/[\ce{A^-}]$}
  \item p$K_a$ of \ce{HC2H3O2}: \answerblank[2cm]{4.74}
  \item $\log(1) =$ \answerblank[1.4cm]{0}
\end{warmup}

\segment{INSTRUCTION A}{25}{Deriving the equation}

\notesectionz{From $K_a$ to a log equation}{15.2}
\practice{5}

Start from the rearranged $K_a$ expression and take $-\log$ of both sides:
\[ [\ce{H+}] = K_a\frac{[\ce{HA}]}{[\ce{A^-}]}
   \;\Longrightarrow\;
   -\log[\ce{H+}] = -\log K_a - \log\frac{[\ce{HA}]}{[\ce{A^-}]} \]
Inverting the fraction flips the sign of its log, which gives the
\term{Henderson--Hasselbalch equation}:
\[ \boxed{\;\mathrm{pH} = \mathrm{p}K_a
   + \log\frac{[\ce{A^-}]}{[\ce{HA}]}
   = \mathrm{p}K_a + \log\frac{[\text{base}]}{[\text{acid}]}\;} \]

Read the equation as three separate statements:

\begin{itemize}[leftmargin=*,itemsep=0.3em]
  \item If $[\text{base}] = [\text{acid}]$, the log term is
        \answerblank[1.4cm]{0}, so the pH equals
        \answerblank[2cm]{p$K_a$}.
  \item More base than acid $\Rightarrow$ log term positive $\Rightarrow$ pH
        \answerblank[2.4cm]{above p$K_a$}.
  \item More acid than base $\Rightarrow$ pH
        \answerblank[2.4cm]{below p$K_a$}.
\end{itemize}

\begin{apctrap}
\textbf{Base over acid, and the base is the anion.} Writing the ratio upside
down flips the sign of the correction, so a pH that should sit above p$K_a$
lands the same distance below it. Before you divide, name the two species
out loud: \ce{C2H3O2^-} is the base, \ce{HC2H3O2} is the acid.
\end{apctrap}

\segment{GUIDED PRACTICE}{15}{Straight substitution}

\begin{enumerate}[leftmargin=*,itemsep=0.4em]
  \item \SI{0.25}{\Molar} \ce{HC2H3O2} and \SI{0.15}{\Molar}
        \ce{NaC2H3O2}: \workspace[1.8cm]
        \keyonly{\begin{apcanswer}pH $= 4.74 + \log(0.15/0.25)
        = 4.74 - 0.22 = \mathbf{4.52}$\end{apcanswer}}
  \item \SI{0.20}{\Molar} \ce{NH3} and \SI{0.30}{\Molar} \ce{NH4Cl}, given
        $K_a(\ce{NH4+}) = 5.6\times10^{-10}$: \workspace[1.8cm]
        \keyonly{\begin{apcanswer}p$K_a = 9.25$;
        pH $= 9.25 + \log(0.20/0.30) = 9.25 - 0.18 =
        \mathbf{9.07}$\end{apcanswer}}
  \item \SI{0.100}{\Molar} \ce{HOCl} and \SI{0.100}{\Molar} \ce{NaOCl}
        ($K_a = 3.5\times10^{-8}$): \answerblank[2.6cm]{pH $= 7.46$}
\end{enumerate}

\begin{apcnote}[On the ammonium p$K_a$]
Two routes give answers differing by $0.01$: $-\log(5.6\times10^{-10}) =
9.25$, while $14.00 - \mathrm{p}K_b = 14.00 - 4.74 = 9.26$. Both are
correct; the gap is rounding in the tabulated constants, nothing more. AP
scoring accepts either. Do not chase the last digit --- pick one route and
be consistent within a problem.
\end{apcnote}

\segment{INSTRUCTION B}{20}{Why the ratio is what matters}

\notesectionz{Dilution does not change a buffer's pH}{15.2}
\practice{5}

Because only the \emph{ratio} appears in the equation, two solutions with
the same ratio have the same pH no matter how concentrated they are:

\begin{center}
\begin{tabular}{@{}lcc@{}}
\hline
\textbf{Solution} & $[\ce{A^-}]/[\ce{HA}]$ & \textbf{pH} \\
\hline
\SI{5.0}{\Molar} \ce{HC2H3O2} $+$ \SI{3.0}{\Molar} \ce{NaC2H3O2}
  & $3.0/5.0 = 0.60$ & 4.52 \\
\SI{0.050}{\Molar} \ce{HC2H3O2} $+$ \SI{0.030}{\Molar} \ce{NaC2H3O2}
  & $0.030/0.050 = 0.60$ & 4.52 \\
\hline
\end{tabular}
\end{center}

Two consequences worth writing down:

\begin{itemize}[leftmargin=*,itemsep=0.3em]
  \item \textbf{Diluting a buffer does not change its pH} (to a good
        approximation) --- both concentrations fall by the same factor, so
        the ratio is \answerblank[2.4cm]{unchanged}. Compare this with a
        weak acid alone, where dilution \emph{does} raise the pH.
  \item \textbf{You may use moles instead of molarity.} Both species share
        the same volume, so the volume cancels out of the ratio. This saves
        a step on nearly every titration problem.
\end{itemize}

\begin{workedexample}[Worked example 2: moles are enough]
What is the pH of a solution made by dissolving \SI{0.100}{\mole} of
\ce{HC2H3O2} and \SI{0.100}{\mole} of \ce{NaC2H3O2} in enough water to make
\SI{500}{\milli\litre}?

There is no need to compute either molarity. The mole ratio is
$0.100/0.100 = 1$, so
$\mathrm{pH} = \mathrm{p}K_a = \mathbf{4.74}$.

The answer would be identical in \SI{2.0}{\litre}.
\end{workedexample}

\segment{APPLICATION}{20}{Working backwards}

\begin{enumerate}[leftmargin=*,itemsep=0.4em]
  \item A buffer made from \ce{HOCl} and \ce{NaOCl} has pH 7.00. Find
        $[\ce{OCl-}]/[\ce{HOCl}]$. \workspace[2.2cm]
        \keyonly{\begin{apcanswer}$7.00 = 7.46 + \log r$, so
        $\log r = -0.46$ and $r = \mathbf{0.35}$. Less base than acid,
        as expected for a pH below p$K_a$.\end{apcanswer}}
  \item A \ce{NH3}/\ce{NH4Cl} buffer is \SI{0.30}{\Molar} in \ce{NH3} and
        \SI{0.15}{\Molar} in \ce{NH4Cl}. Find the pH and say whether it lies
        above or below p$K_a$, with a reason. \workspace[2.4cm]
        \keyonly{\begin{apcanswer}pH $= 9.25 + \log(0.30/0.15)
        = 9.25 + 0.30 = \mathbf{9.55}$. Above p$K_a$, because the base
        outnumbers its conjugate acid.\end{apcanswer}}
\end{enumerate}

\begin{exitticket}
Two beakers hold acetate buffers. Beaker A is \SI{1.0}{\Molar} in each
component; beaker B is \SI{0.010}{\Molar} in each. Which has the higher pH,
and which is the better buffer?
\keyonly{\begin{apcanswer}Neither has a higher pH --- both ratios are 1, so
both are at pH 4.74. Beaker A is the better buffer, because it holds far
more of each component and can therefore absorb far more added acid or base
before the ratio moves appreciably.\end{apcanswer}}
\end{exitticket}

% =====================================================================
\blocklesson{Buffer Action and Buffering Capacity}{Zumdahl \S15.2--15.3}

\begin{objectives}
  \item Calculate the pH of a buffer after a strong acid or base is added.
  \item Compare buffering capacity and choose components for a target pH.
\end{objectives}

\reading{Zumdahl \S15.2--15.3}{761--770}

\begin{warmup}[3]
  \item Henderson--Hasselbalch: pH $=$
        \answerblank[5cm]{p$K_a + \log([\text{base}]/[\text{acid}])$}
  \item pH of a buffer with equal amounts of both components:
        \answerblank[2cm]{p$K_a$}
  \item Diluting a buffer changes its pH how?
        \answerblank[3cm]{hardly at all}
\end{warmup}

\segment{INSTRUCTION A}{25}{Stoichiometry first, then equilibrium}

\notesectionz{The two-step method}{15.2}
\practice{5}

This is the most important procedure in the chapter, and the one students
most often try to skip.

\begin{enumerate}[leftmargin=*,itemsep=0.35em]
  \item \textbf{Stoichiometry step.} Strong acid or strong base reacts
        \term{completely}. Added \ce{OH-} converts \ce{HA} to \ce{A-};
        added \ce{H+} converts \ce{A-} to \ce{HA}. Work in
        \answerblank[2cm]{moles} (or mmol) and take the reaction to
        completion. No $K$ appears in this step.
  \item \textbf{Equilibrium step.} Take the new amounts and put them into
        Henderson--Hasselbalch.
\end{enumerate}

\begin{apctrap}
Do \textbf{not} put the added \ce{OH-} into an ICE table. Strong base does
not sit at equilibrium with a weak acid --- it consumes it. Mixing the two
steps is the single most common way to lose every point on a buffer
question.
\end{apctrap}

\begin{workedexample}[Worked example 3: a buffer absorbing a strong base]
\SI{1.00}{\litre} of solution is \SI{0.50}{\Molar} in \ce{HC2H3O2} and
\SI{0.50}{\Molar} in \ce{NaC2H3O2}. Add \SI{0.010}{\mole} of \ce{NaOH}.

\textbf{Start:} ratio $=1$, so pH $= 4.74$.

\textbf{Step 1 --- stoichiometry.} $\ce{OH^- + HC2H3O2 -> C2H3O2^- + H2O}$

\begin{center}
\begin{tabular}{@{}lccc@{}}
\hline
 & \ce{HC2H3O2} & \ce{OH-} & \ce{C2H3O2^-} \\
\hline
before & 0.50 & 0.010 & 0.50 \\
after  & 0.49 & 0      & 0.51 \\
\hline
\end{tabular}
\end{center}

\textbf{Step 2 --- equilibrium.}
$\mathrm{pH} = 4.74 + \log(0.51/0.49) = 4.74 + 0.02 = \mathbf{4.76}$

\medskip
\textbf{For contrast:} the same \SI{0.010}{\mole} of \ce{NaOH} in
\SI{1.00}{\litre} of \emph{pure water} gives $[\ce{OH-}] = 0.010$,
pOH $= 2.00$, and $\mathrm{pH} = \mathbf{12.00}$.

The buffer moved \textbf{0.02} pH units. Water moved \textbf{5.00}.
\end{workedexample}

\segment{GUIDED PRACTICE}{15}{The same buffer, acid added}

Add \SI{0.010}{\mole} of \ce{HCl} to the same original buffer instead.

\begin{enumerate}[leftmargin=*,itemsep=0.35em]
  \item Which component is consumed? \answerblank[3.4cm]{\ce{C2H3O2^-}}
  \item New amounts: \ce{HC2H3O2} $=$ \answerblank[1.6cm]{0.51},
        \ce{C2H3O2^-} $=$ \answerblank[1.6cm]{0.49}
  \item New pH: \answerblank[2.4cm]{4.72}
  \item Same \ce{HCl} in \SI{1.00}{\litre} pure water:
        \answerblank[2.4cm]{2.00}
\end{enumerate}

\segment{INSTRUCTION B}{20}{Capacity, and choosing a buffer}

\notesectionz{Why equal amounts buffer best}{15.3}
\practice{6}

\term{Buffering capacity} is how much added acid or base a buffer can absorb
before its pH moves appreciably. Two things control it:

\begin{itemize}[leftmargin=*,itemsep=0.3em]
  \item \textbf{How much} of each component is present. More is better ---
        this is why dilution, which leaves the pH alone, still \emph{hurts}
        the buffer.
  \item \textbf{How lopsided} the ratio is. Zumdahl's Table 15.1 makes the
        case with two solutions at very different ratios, each receiving
        \SI{0.01}{\mole} of \ce{H+}:
\end{itemize}

\begin{center}
\begin{tabular}{@{}lcccc@{}}
\hline
 & $[\ce{A^-}]/[\ce{HA}]$ before & after & change & pH shift \\
\hline
A \ \SI{1.00}{\Molar} / \SI{1.00}{\Molar} & 1.00 & 0.98 & 2.0\% &
  $4.74 \to 4.73$ \\
B \ \SI{1.00}{\Molar} / \SI{0.01}{\Molar} & 100  & 49.5 & 50.5\% &
  $6.74 \to 6.43$ \\
\hline
\end{tabular}
\end{center}

Same amount of acid added; solution B's ratio changed
\answerblank[2.6cm]{25 times} as much. The lesson:

\begin{center}
\fbox{\parbox{0.86\linewidth}{\centering
Optimal buffering occurs when $[\ce{HA}] = [\ce{A^-}]$, that is, when
$\mathrm{pH} = \mathrm{p}K_a$. \\
A buffer is useful over roughly
$\mathrm{p}K_a \pm 1$.}}
\end{center}

So to \textbf{design} a buffer for a target pH, choose a weak acid whose
p$K_a$ is as close as possible to that pH, then set the ratio with
Henderson--Hasselbalch.

\begin{workedexample}[Worked example 4: designing a pH 4.00 buffer]
Which acid, and in what ratio?

\begin{center}
\begin{tabular}{@{}lcc@{}}
\hline
\textbf{Acid} & p$K_a$ & $|\mathrm{pH} - \mathrm{p}K_a|$ \\
\hline
\ce{HF}       & 3.14 & 0.86 \\
\ce{HNO2}     & 3.40 & \textbf{0.60} \\
\ce{HC2H3O2}  & 4.74 & 0.74 \\
\ce{HOCl}     & 7.46 & 3.46 \quad (unusable) \\
\hline
\end{tabular}
\end{center}

\ce{HNO2} is the closest, so use \ce{HNO2}/\ce{NaNO2}:
\[ 4.00 = 3.40 + \log\frac{[\ce{NO2^-}]}{[\ce{HNO2}]}
   \;\Longrightarrow\;
   \frac{[\ce{NO2^-}]}{[\ce{HNO2}]} = 10^{0.60} = \mathbf{4.0} \]
Four times as much nitrite as nitrous acid. \ce{HOCl} is useless here: its
p$K_a$ is 3.5 units away, so reaching pH 4.00 would need a ratio near
$10^{-3.5}$ --- almost no conjugate base present, and no capacity against
added acid.
\end{workedexample}

\segment{APPLICATION}{20}{Capacity reasoning}

\begin{enumerate}[leftmargin=*,itemsep=0.4em]
  \item \SI{1.0}{\litre} of a buffer is \SI{0.40}{\Molar} in
        \ce{HC2H3O2} and \SI{0.40}{\Molar} in \ce{NaC2H3O2}. Add
        \SI{0.050}{\mole} \ce{NaOH}. Find the new pH. \workspace[2.6cm]
        \keyonly{\begin{apcanswer}After reaction: \ce{HC2H3O2} $= 0.35$,
        \ce{C2H3O2^-} $= 0.45$.
        pH $= 4.74 + \log(0.45/0.35) = 4.74 + 0.11 =
        \mathbf{4.85}$\end{apcanswer}}
  \item Explain why a buffer eventually fails if you keep adding base.
        \answerlines[3]{Each addition converts more \ce{HA} into \ce{A-}.
        Once the \ce{HA} is essentially gone there is nothing left to
        neutralize the base, so further additions simply raise
        $[\ce{OH-}]$ and the pH climbs sharply.}
  \item Rank for buffering capacity at pH 4.74, highest first:
        (i) \SI{1.0}{\Molar}/\SI{1.0}{\Molar} acetate,
        (ii) \SI{0.10}{\Molar}/\SI{0.10}{\Molar} acetate,
        (iii) \SI{1.0}{\Molar}/\SI{0.010}{\Molar} acetate.
        \answerlines[2]{(i) $>$ (ii) $>$ (iii). (i) beats (ii) on amount at
        the same ideal ratio; (iii) is at neither the right pH nor a
        balanced ratio, and has almost no acid to absorb added base.}
\end{enumerate}

\begin{exitticket}
Blood is buffered near pH 7.4 by the \ce{H2CO3}/\ce{HCO3-} pair
(p$K_a \approx 6.4$). Is this the ideal buffer for that pH? Why might the
body use it anyway?
\keyonly{\begin{apcanswer}Not ideal --- 7.4 is a full unit above p$K_a$, at
the edge of the useful range, so the ratio is about 10:1 and the capacity
against added base is limited. The body uses it because \ce{CO2} is
continuously produced and can be exhaled, so the system is constantly
replenished and is not a closed buffer at all.\end{apcanswer}}
\end{exitticket}

% =====================================================================
\blocklesson{Titrations and pH Curves}{Zumdahl \S15.4}

\begin{objectives}
  \item Calculate pH at any point of a strong--strong or weak--strong
        titration.
  \item Interpret the regions and landmarks of a titration curve.
\end{objectives}

\reading{Zumdahl \S15.4}{770--785}

\begin{warmup}[4]
  \item Buffer method, step 1: \answerblank[5.6cm]{stoichiometry to
        completion}
  \item Buffer method, step 2: \answerblank[4.6cm]{Henderson--Hasselbalch}
  \item mmol $=$ \answerblank[3.6cm]{mL $\times$ molarity}
  \item pH of a \SI{0.10}{\Molar} \ce{NaC2H3O2} solution:
        \answerblank[2cm]{8.87}
\end{warmup}

\segment{INSTRUCTION A}{25}{Strong acid with strong base}

\notesectionz{Bookkeeping in millimoles}{15.4}
\practice{5}

Since $\text{mL}\times\text{molarity} = \text{mmol}$, working in mL and mmol
avoids every unit conversion. The total volume is always
\answerblank[6.2cm]{initial volume $+$ volume added}.

\begin{workedexample}[Worked example 5: Zumdahl's nitric acid titration]
\SI{50.0}{\milli\litre} of \SI{0.200}{\Molar} \ce{HNO3}
($= \SI{10.0}{\milli\mole}$ \ce{H+}) titrated with \SI{0.100}{\Molar}
\ce{NaOH}. The reaction is $\ce{H+ + OH^- -> H2O}$.

\begin{center}
\begin{tabular}{@{}rrrr@{}}
\hline
\textbf{mL NaOH} & \textbf{mmol \ce{OH-}} & \textbf{excess} &
  \textbf{pH} \\
\hline
0     & 0    & 10.0 mmol \ce{H+} & 0.70 \\
10.0  & 1.00 & 9.00 mmol \ce{H+} & 0.82 \\
20.0  & 2.00 & 8.00 mmol \ce{H+} & 0.94 \\
50.0  & 5.00 & 5.00 mmol \ce{H+} & 1.30 \\
100.0 & 10.0 & none --- \textbf{equivalence} & \textbf{7.00} \\
150.0 & 15.0 & 5.00 mmol \ce{OH-} & 12.40 \\
\hline
\end{tabular}
\end{center}

Sample: at \SI{20.0}{\milli\litre},
$[\ce{H+}] = 8.00/(50.0+20.0) = \SI{0.114}{\Molar}$, pH $= 0.94$.

At equivalence the only species left are \ce{Na+}, \ce{NO3-} and water.
Neither ion reacts, so the pH is exactly \textbf{7.00}.
\end{workedexample}

\segment{INSTRUCTION B}{20}{Weak acid with strong base}

\notesectionz{Four regions, four different methods}{15.4}
\practice{5}

This is the curve AP asks about most, and the reason is that
\emph{each region needs a different calculation}:

\begin{center}
\begin{tabular}{@{}llp{6.4cm}@{}}
\hline
\textbf{Region} & \textbf{Major species} & \textbf{Method} \\
\hline
Before any base & \ce{HA} only & weak-acid ICE, $x=\sqrt{K_aC_0}$ \\
Before equivalence & \ce{HA} $+$ \ce{A-} & \textbf{buffer:}
  Henderson--Hasselbalch \\
At equivalence & \ce{A-} only & weak-\emph{base} ICE, $K_b = K_w/K_a$ \\
After equivalence & \ce{A-} $+$ excess \ce{OH-} & excess strong base alone \\
\hline
\end{tabular}
\end{center}

\begin{workedexample}[Worked example 6: acetic acid titrated with \ce{NaOH}]
\SI{50.0}{\milli\litre} of \SI{0.10}{\Molar} \ce{HC2H3O2}
($= \SI{5.0}{\milli\mole}$) with \SI{0.10}{\Molar} \ce{NaOH}. Equivalence
comes at \SI{50.0}{\milli\litre}.

\begin{center}
\begin{tabular}{@{}rllr@{}}
\hline
\textbf{mL} & \textbf{Region} & \textbf{Work} & \textbf{pH} \\
\hline
0    & weak acid & $x=\sqrt{(1.8\times10^{-5})(0.10)}=1.34\times10^{-3}$
  & 2.87 \\
10.0 & buffer & $4.74+\log(1.0/4.0)$ & 4.14 \\
25.0 & \textbf{halfway} & $4.74+\log(2.5/2.5)$ & \textbf{4.74} \\
40.0 & buffer & $4.74+\log(4.0/1.0)$ & 5.34 \\
50.0 & \textbf{equivalence} & $[\ce{A^-}]=0.050$, $K_b=5.6\times10^{-10}$
  & \textbf{8.72} \\
60.0 & excess base & $[\ce{OH-}]=1.0/110.0$ & 11.96 \\
\hline
\end{tabular}
\end{center}

\textbf{The equivalence-point calculation in full.} All \SI{5.0}{\milli\mole}
of acid has become acetate, in $50.0+50.0 = \SI{100.0}{\milli\litre}$:
$[\ce{C2H3O2^-}] = \SI{0.050}{\Molar}$. Then
$x = \sqrt{(5.6\times10^{-10})(0.050)} = 5.3\times10^{-6} = [\ce{OH-}]$,
pOH $= 5.28$, pH $= \mathbf{8.72}$.
\end{workedexample}

\begin{apctrap}
\textbf{The equivalence point of a weak acid titration is not pH 7.} All the
acid has been converted to its conjugate base, and that base hydrolyses. For
a weak acid titrated with a strong base the equivalence pH is
\answerblank[1.8cm]{above 7}; for a weak base titrated with a strong acid it
is \answerblank[1.8cm]{below 7}. Only strong-with-strong lands on 7.00.
\end{apctrap}

\begin{apcnote}[Two useful landmarks]
\textbf{Halfway point.} Exactly half the acid has been converted, so
$[\ce{HA}]=[\ce{A^-}]$ and $\mathrm{pH} = \mathrm{p}K_a$. This is how $K_a$
is measured from a curve, and it is worth a point almost every time it
appears.

\textbf{Equivalence point.} Moles of added base $=$ moles of acid
originally present. This is a \emph{stoichiometric} definition --- it says
nothing about the pH being 7.
\end{apcnote}

\begin{center}
\begin{tikzpicture}
\begin{axis}[width=11cm, height=6.4cm,
    xlabel={mL of \SI{0.100}{\Molar} \ce{NaOH} added},
    ylabel={pH}, xmin=0, xmax=80, ymin=0, ymax=13.5,
    xtick={0,20,40,50,60,80}, ytick={0,2,4,6,8,10,12},
    axis lines=left, legend style={draw=none, font=\sffamily\scriptsize,
      at={(0.98,0.02)}, anchor=south east}]
  \addplot[seriesA] coordinates {(0,1.000) (5,1.087) (10,1.176) (20,1.368)
    (30,1.602) (40,1.954) (45,2.279) (48,2.690) (49,2.996) (49.5,3.299)
    (49.9,4.000) (50,7.000) (50.1,10.000) (50.5,10.697) (51,10.996)
    (52,11.292) (55,11.678) (60,11.959) (70,12.222) (80,12.363)};
  \addlegendentry{\SI{0.100}{\Molar} \ce{HCl} (strong)}
  \addplot[seriesB] coordinates {(0,2.875) (5,3.799) (10,4.145) (20,4.569)
    (30,4.921) (40,5.347) (45,5.699) (48,6.125) (49,6.435) (49.5,6.740)
    (49.9,7.442) (50,8.722) (50.1,10.001) (50.5,10.697) (51,10.996)
    (52,11.292) (55,11.678) (60,11.959) (70,12.222) (80,12.363)};
  \addlegendentry{\SI{0.100}{\Molar} \ce{HC2H3O2} (weak)}
  \draw[densely dotted] (axis cs:25,0) -- (axis cs:25,4.74);
  \draw[densely dotted] (axis cs:0,4.74) -- (axis cs:25,4.74);
  \node[font=\scriptsize,anchor=west] at (axis cs:26,4.0)
    {halfway: pH $=$ p$K_a$ $= 4.74$};
  % mark both equivalence points explicitly -- they share a volume but not a pH
  \fill (axis cs:50,7.00) circle (1.6pt);
  \fill (axis cs:50,8.72) circle (1.6pt);
  \node[font=\scriptsize,anchor=east] at (axis cs:47.5,8.9)
    {equiv.\ 8.72 (weak)};
  \node[font=\scriptsize,anchor=west] at (axis cs:52,6.4)
    {equiv.\ 7.00 (strong)};
\end{axis}
\end{tikzpicture}
\end{center}

Both curves start and end in different places but converge after
equivalence, because once the acid is gone the pH is set only by
\answerblank[3.6cm]{excess \ce{NaOH}}, which is the same in both flasks.

\segment{APPLICATION}{20}{Reading and calculating}

\begin{enumerate}[leftmargin=*,itemsep=0.4em]
  \item \SI{100.0}{\milli\litre} of \SI{0.050}{\Molar} \ce{NH3}
        ($K_b = 1.8\times10^{-5}$) is titrated with \SI{0.10}{\Molar}
        \ce{HCl}. Find the volume at equivalence and the pH there.
        \workspace[3cm]
        \keyonly{\begin{apcanswer}\SI{5.0}{\milli\mole} \ce{NH3} needs
        \SI{5.0}{\milli\mole} \ce{HCl} $=$ \SI{50.0}{\milli\litre}.
        At equivalence $[\ce{NH4+}] = 5.0/150.0 = 0.0333$;
        $K_a = 5.6\times10^{-10}$;
        $x = \sqrt{(5.6\times10^{-10})(0.0333)} = 4.3\times10^{-6}$;
        pH $= \mathbf{5.36}$ --- \textbf{acidic}, as expected for the
        conjugate acid of a weak base.\end{apcanswer}}
  \item For that same titration, what is the pH after
        \SI{25.0}{\milli\litre} of \ce{HCl}, and why can you write it down
        with almost no work? \answerlines[2]{That is the halfway point, so
        $[\ce{NH3}] = [\ce{NH4+}]$ and pH $=$ p$K_a(\ce{NH4+}) = 9.25$.}
  \item A curve for a monoprotic acid shows its equivalence point at pH
        9.1. Is the acid strong or weak? Justify.
        \answerlines[2]{Weak. A strong acid titrated with strong base gives
        equivalence at exactly 7.00; a basic equivalence point means the
        anion left behind hydrolyses, which only happens if it is the
        conjugate base of a weak acid.}
\end{enumerate}

\begin{exitticket}
Why does the buffer region of the weak-acid curve appear flat, while the
strong-acid curve has no flat region at all?
\keyonly{\begin{apcanswer}In the weak-acid case both \ce{HA} and \ce{A-} are
present, so added \ce{OH-} only changes their ratio --- and pH depends on
the \emph{log} of that ratio, so it moves slowly. In the strong-acid case
there is no conjugate pair to absorb the base; every added \ce{OH-}
directly removes \ce{H+} from solution.\end{apcanswer}}
\end{exitticket}

% =====================================================================
\blocklesson{Indicators and Titration as Analysis}{Zumdahl \S15.5--15.6}

\begin{objectives}
  \item Select an indicator whose range matches a titration's equivalence
        point.
  \item Determine $K_a$ and concentration from titration data.
\end{objectives}

\reading{Zumdahl \S15.5--15.6}{785--792}

\begin{warmup}[3]
  \item Equivalence pH, weak acid $+$ strong base:
        \answerblank[2.2cm]{above 7}
  \item Equivalence pH, weak base $+$ strong acid:
        \answerblank[2.2cm]{below 7}
  \item At the halfway point, pH $=$ \answerblank[2cm]{p$K_a$}
\end{warmup}

\segment{INSTRUCTION A}{25}{How an indicator works}

\notesectionz{An indicator is just another weak acid}{15.5}
\practice{6}

An \term{acid--base indicator} is a weak acid, written \ce{HIn}, whose two
forms have different colors:
\[ \ce{HIn(aq) <=> H+(aq) + In^-(aq)} \]
\centerline{\small (one color) \hfill (the other color)}

Because it is a weak acid, it obeys the same equation everything else in
this chapter obeys:
\[ \mathrm{pH} = \mathrm{p}K_a + \log\frac{[\ce{In^-}]}{[\ce{HIn}]} \]

The eye needs about \textbf{one part in ten} of the new form before it
registers a change in color. So, titrating an \emph{acid} (going from
\ce{HIn} toward \ce{In-}), the change first shows when
$[\ce{In^-}]/[\ce{HIn}] = 1/10$:
\[ \mathrm{pH} = \mathrm{p}K_a + \log\tfrac{1}{10}
   = \mathrm{p}K_a - 1 \]
and titrating a \emph{base}, at the reciprocal ratio, at
$\mathrm{p}K_a + 1$. Hence the general rule:

\begin{center}
\fbox{\parbox{0.78\linewidth}{\centering An indicator changes color over
roughly $\mathrm{p}K_a \pm 1$ --- a span of about \textbf{2 pH units}.}}
\end{center}

\begin{workedexample}[Worked example 7: bromthymol blue]
Bromthymol blue has $K_a = 1.0\times10^{-7}$, yellow as \ce{HIn} and blue as
\ce{In-}. An acidic solution containing it is titrated with \ce{NaOH}. At
what pH does the color first visibly change?

$\mathrm{p}K_a = 7.00$, and the change is visible at
$[\ce{In^-}]/[\ce{HIn}] = 1/10$:
\[ \mathrm{pH} = 7.00 + \log(0.10) = 7.00 - 1.00 = \mathbf{6.00} \]
A greenish tint --- a little blue mixed into the yellow. Its useful range
runs from about \textbf{6 to 8}.
\end{workedexample}

\notesectionz{End point is not equivalence point}{15.5}

\begin{itemize}[leftmargin=*,itemsep=0.3em]
  \item The \term{equivalence point} is defined by
        \answerblank[4.6cm]{reaction stoichiometry}.
  \item The \term{end point} is defined by
        \answerblank[5.6cm]{the indicator changing color}.
\end{itemize}

They are different things. A good indicator makes them nearly coincide; a
badly chosen one does not.

\segment{GUIDED PRACTICE}{15}{Choosing an indicator}

\begin{center}
\begin{tabular}{@{}lcc@{}}
\hline
\textbf{Indicator} & \textbf{approx.\ p$K_a$} & \textbf{useful range} \\
\hline
Methyl orange     & 3.7 & 3.1--4.4 \\
Methyl red        & 5.3 & 4.4--6.2 \\
Bromthymol blue   & 7.0 & 6.0--7.6 \\
Phenolphthalein   & 9.1 & 8.2--10.0 \\
\hline
\end{tabular}
\end{center}

\begin{enumerate}[leftmargin=*,itemsep=0.35em]
  \item \ce{HCl} with \ce{NaOH}, equivalence pH 7.00. Suitable indicators:
        \answerblank[8.5cm]{any of them --- the jump is nearly vertical}
  \item \ce{HC2H3O2} with \ce{NaOH}, equivalence pH 8.72:
        \answerblank[6cm]{phenolphthalein}
  \item Why is methyl red wrong for item 2?
        \answerlines[2]{Its range (4.4--6.2) sits inside the buffer region,
        well before equivalence. It would change color while acid still
        remained, so the end point would badly underestimate the volume
        needed.}
  \item \ce{NH3} with \ce{HCl}, equivalence pH 5.36:
        \answerblank[6cm]{methyl red}
\end{enumerate}

\segment{INSTRUCTION B}{20}{Titration as a measuring tool}

\notesectionz{Getting $K_a$ and concentration from a curve}{15.4}
\practice{4}

A titration curve carries two independent pieces of information:

\begin{itemize}[leftmargin=*,itemsep=0.3em]
  \item The \textbf{equivalence volume} gives the amount --- and therefore
        the concentration --- of the unknown.
  \item The \textbf{pH at half that volume} gives
        \answerblank[2cm]{p$K_a$}, and hence $K_a$.
\end{itemize}

\begin{workedexample}[Worked example 8: identifying an unknown acid]
\SI{25.00}{\milli\litre} of a monoprotic acid requires
\SI{18.50}{\milli\litre} of \SI{0.1050}{\Molar} \ce{NaOH} to reach
equivalence. The pH at \SI{9.25}{\milli\litre} is 4.20.

\textbf{Concentration.} $n = (18.50)(0.1050) = \SI{1.943}{\milli\mole}$, so
$C = 1.943/25.00 = \SI{0.0777}{\Molar}$.

\textbf{$K_a$.} \SI{9.25}{\milli\litre} is exactly half of
\SI{18.50}{\milli\litre}, so this is the halfway point:
$\mathrm{p}K_a = 4.20$ and
$K_a = 10^{-4.20} = \mathbf{6.3\times10^{-5}}$.

\textbf{Sanity check.} A $K_a$ near $10^{-5}$ is a typical weak organic
acid, and it predicts an equivalence point above 7 --- consistent with
titrating a weak acid with a strong base.
\end{workedexample}

\notesectionz{Polyprotic curves}{15.6}

\begin{apcnote}[Enrichment --- shape only]
A polyprotic acid is titrated one proton at a time, giving \emph{one
equivalence point per ionizable proton}. Each addition of \ce{OH-} strips
off the next proton in turn:
\[ \ce{H3PO4 -> H2PO4^- -> HPO4^2- -> PO4^3-} \]
Each step consumes the \emph{same} number of moles of base, so the
equivalence points are evenly spaced along the volume axis --- which is how
you count the protons from a curve. Between them sit buffer plateaus.

You should be able to read a polyprotic curve and say how many acidic
protons the acid has. Calculating pH at the second and third equivalence
points is beyond the current CED.
\end{apcnote}

\segment{APPLICATION}{20}{Putting it together}

\begin{enumerate}[leftmargin=*,itemsep=0.4em]
  \item A \SI{25.0}{\milli\litre} sample of \SI{0.100}{\Molar}
        \ce{HC2H3O2} is titrated with \SI{0.100}{\Molar} \ce{NaOH}. After
        \SI{10.0}{\milli\litre} the pH is measured. Predict it.
        \workspace[2.4cm]
        \keyonly{\begin{apcanswer}Start: \SI{2.50}{\milli\mole} \ce{HA}.
        Added: \SI{1.00}{\milli\mole} \ce{OH-}. Left: 1.50 \ce{HA} and
        1.00 \ce{A-}.
        pH $= 4.74 + \log(1.00/1.50) = 4.74 - 0.18 =
        \mathbf{4.56}$\end{apcanswer}}
  \item An indicator has $K_a = 1\times10^{-5}$. Over what pH range does it
        change color, and which titration in this block would it suit?
        \answerlines[2]{p$K_a = 5.0$, so the range is about 4 to 6. It suits
        the \ce{NH3}/\ce{HCl} titration (equivalence 5.36), not the acetic
        acid titration (equivalence 8.72).}
  \item A student titrating acetic acid with \ce{NaOH} uses methyl orange
        and reports a concentration that is far too low. Explain the
        source of the error. \answerlines[3]{Methyl orange changes color
        near pH 3--4, which occurs very early in this titration while most
        of the acetic acid is still unreacted. The student stopped adding
        base long before equivalence, recorded too small a volume, and so
        calculated too few moles of acid.}
\end{enumerate}

\begin{exitticket}
State the one question you should ask before starting any titration
calculation, and why it settles the method.
\keyonly{\begin{apcanswer}``What are the major species after the strong
acid or base has reacted completely?'' The answer tells you which of the
four situations you are in --- weak acid alone, buffer, conjugate base
alone, or excess strong reagent --- and each has its own method. Choosing
the method first and identifying the species afterwards is how students end
up applying Henderson--Hasselbalch at an equivalence point, where there is
no buffer at all.\end{apcanswer}}
\end{exitticket}

\end{document}
